\section*{Declaration of competing interest}
The author declares no competing financial interests or personal relationships that could have appeared to
influence the work reported in this paper.

\section*{Funding}
This research did not receive any specific grant from funding agencies in the public, commercial, or not-for-profit
sectors.

\section*{Data and code availability}
The code, evaluation texts, per-utterance scores and the scripts that produce every number, table and figure are
archived with the paper sources on Zenodo (\zenodo); the engine source also ships with the demo. The model is at \hfmodel and the demo at \hfspace. FLEURS \citep{conneau2023fleurs}, lahgtna-v2 \citep{oddadmix2026lahgtnadata}
and NileTTS \citep{khamis2026niletts} are public datasets; FLEURS reference clips and the audio cloned from them are not
redistributed (\cref{sec:release}).

\section*{Ethics}
\emph{Listening test.} Native Egyptian adults took part in the blind listening test (\cref{sec:dissociation},
\cref{app:listening}) as unpaid volunteers. The first page of the test stated the purpose of the study, that taking part was
voluntary and could be stopped at any time by closing the page, what would be stored, and that the results would be
published only in aggregate, and it gave a contact address; listeners confirmed that they were 18 or older and agreed to
take part before the first clip. The test asked for no name or contact details and its database stored no IP address:
each session holds a random identifier, the ratings and four coarse descriptors (age band, gender, governorate,
listening device). Individual answers are not published. Listeners heard speech cloned from FLEURS speakers, released for
speech research, and real clips from the public lahgtna-v2 corpus; they were asked not to record or share the clips, and
no new recordings of people were made. No institutional ethics review was sought: taking part was anonymous, voluntary and unpaid, consisted only of listening to short clips and rating them, carried no foreseeable risk beyond those of everyday life, and collected no information that identifies a person.

\emph{Voice cloning.} Voice cloning can impersonate people. The safeguards of the release (a consent statement in the
demo, no storage of uploads, an AudioSeal watermark on every output) and their limits are described in
\cref{sec:release}.

\section*{Declaration of generative AI and AI-assisted technologies in the manuscript preparation process}
During the preparation of this work, the author used Claude (Anthropic) for help with code and with the academic
writing. After using this tool, the author reviewed and edited the content as needed and takes full responsibility for
the content of the published article.
